% GENERATED FILE. Edit canonical lessons, then run npm run generate:books.
% canonical-source-hash: fnv1a64:dee639945cd9c1e2
% canonical-lessons: MR-C46-huun, MR-C46-paasun, MR-C46-kuthun, MR-R46-from

\chapter{Where You Came From}
\label{ch:mr-from}
\hlchaptermodality{46}

\begin{chapteropening}
\textbf{By the end of this chapter:} \emph{I can say where I come from and ask somebody else, and I can start a stretch of time as well as a journey.}

Complete MR-R46-from and run the interview's second exchange both ways -- ask kuṭhūn yetā?, answer puṇyāhūn -- then start a stretch of time with udyāpāsūn.
\end{chapteropening}

\section[-hūn]{-\mr{हून} (-hūn) — from}

\label{lesson:MR-C46-huun}



\begin{quote}
Chapter twenty-three taught \textbf{\mr{शहर}} as a line on a form. A form wants the name of a city. An interviewer wants something else.

\begin{itemize}
\raggedright
  \item \emph{Recall:} say \emph{-ne}, then \emph{-śī}, and name which one takes a person
\end{itemize}
\end{quote}

\subsection*{You'll want to know: -\mr{हून}}
\begin{quote}
\textbf{-\mr{हून}} — \emph{-hūn} — \textbf{from}
\end{quote}

\begin{quote}
\textbf{\mr{मी} \mr{पुण्याहून} \mr{येतो}.} — \emph{mī puṇyāhūn yeto.} — \textbf{I come from Pune.}
\end{quote}

That sentence is the second thing said in almost every spoken A1 interview, and until this lesson the book could write \textbf{\mr{शहर}: \mr{पुणे}} on a form and not say it out loud. A form field is a label; \textbf{-\mr{हून}} is an answer.

The stem is the one you already bend: \textbf{\mr{पुणे} $\to$ \mr{पुण्या}-}, then \textbf{-\mr{हून}}, no space. And because \textbf{\mr{इथे}} and \textbf{\mr{तिथे}} are place words too, they take it and give you two more:

\begin{quote}
\textbf{\mr{इथून}} — \emph{ithūn} — \textbf{from here}
\end{quote}

\begin{quote}
\textbf{\mr{तिथून}} — \emph{tithūn} — \textbf{from there}
\end{quote}

Notice what happened to the shape: \textbf{\mr{इथे}} ends in \textbf{-\mr{ए}} and the ending replaced it, exactly as it does on \textbf{\mr{पुणे}}. The rule you learned in chapter forty-one is doing all the work.

\subsection*{Your turn}
Say these aloud:

\begin{itemize}
\raggedright
  \item \emph{mī puṇyāhūn yeto} — I come from Pune
  \item \emph{ithūn}, then \emph{tithūn} — from here, from there
  \item \emph{gharāhūn} — from the house
\end{itemize}

\begin{tcolorbox}[breakable,colback=teal!4,colframe=teal!35!black,title={Before you move on}]
What can you now say that a form field could not? (\emph{\textbf{I come from Pune}} — an answer, not a label.) What does \textbf{\mr{इथे}} become with it? (\textbf{\mr{इथून}}.)

Sources: \href{https://en.wiktionary.org/wiki/\%E0\%A4\%B9\%E0\%A5\%82\%E0\%A4\%A8}{Wiktionary: \mr{हून}}.
\end{tcolorbox}
\section[pāsūn]{\mr{पासून} (pāsūn) — from, since}

\label{lesson:MR-C46-paasun}



\begin{quote}
\textbf{-\mr{हून}} says where a journey began. It does not say when one did.
\end{quote}

\subsection*{You'll want to know: \mr{पासून}}
\begin{quote}
\textbf{\mr{पासून}} — \emph{pāsūn} — \textbf{from, since}
\end{quote}

\begin{quote}
\textbf{\mr{उद्यापासून}} — \emph{udyāpāsūn} — \textbf{from tomorrow onwards}
\end{quote}

\textbf{\mr{पासून}} is written joined, like \textbf{-\mr{हून}}, and takes the same oblique stem. What makes it worth its own lesson is that it works on \textbf{time} as well as place — Marathi uses one word where English splits \emph{from} and \emph{since}, and a Marathi speaker feels no difference between the start of a road and the start of a week.

Against \textbf{-\mr{हून}} the contrast is direction of attention. \textbf{-\mr{हून}} looks back at the place you LEFT; \textbf{\mr{पासून}} looks forward from a starting point that is still running:

\begin{quote}
\textbf{\mr{पुण्याहून}} — I came \emph{from} Pune, and I am elsewhere now.
\end{quote}

\begin{quote}
\textbf{\mr{उद्यापासून}} — \emph{starting} tomorrow, and continuing.
\end{quote}

That is why a notice about opening hours uses \textbf{\mr{पासून}} and a traveller uses \textbf{-\mr{हून}}.

\subsection*{Your turn}
Say these aloud:

\begin{itemize}
\raggedright
  \item \emph{udyāpāsūn} — from tomorrow
  \item \emph{gharāpāsūn} — from the house onwards
  \item the pair — \emph{puṇyāhūn}, \emph{udyāpāsūn}
\end{itemize}

\begin{tcolorbox}[breakable,colback=teal!4,colframe=teal!35!black,title={Before you move on}]
Which of the two also does TIME? (\textbf{\mr{पासून}}.) Which one looks back at where you left? (\textbf{-\mr{हून}}.)

Sources: \href{https://en.wiktionary.org/wiki/\%E0\%A4\%AA\%E0\%A4\%BE\%E0\%A4\%B8\%E0\%A5\%82\%E0\%A4\%A8}{Wiktionary: \mr{पासून}}.
\end{tcolorbox}
\section[kuṭhūn]{\mr{कुठून} (kuṭhūn) — from where}

\label{lesson:MR-C46-kuthun}



\begin{quote}
You can answer \emph{I come from Pune}. Nobody has asked you yet.
\end{quote}

\subsection*{You'll want to know: \mr{कुठून}}
\begin{quote}
\textbf{\mr{कुठून}} — \emph{kuṭhūn} — \textbf{from where}
\end{quote}

\begin{quote}
\textbf{\mr{तुम्ही} \mr{कुठून} \mr{येता}?} — \emph{tumhī kuṭhūn yetā?} — \textbf{where are you from?}
\end{quote}

Nothing here is new except the answer to a question you can now finally ask. \textbf{\mr{कुठून}} is \textbf{\mr{कुठे}} with \textbf{-\mr{हून}} on it, and the \textbf{-\mr{ए}} dropped for \textbf{-\mr{ऊ}} in exactly the way \textbf{\mr{इथे}} became \textbf{\mr{इथून}} two lessons ago. A place word, an ablative, and the same rule.

Its position obeys chapter thirty-eight: an asking word stands in the answer's own slot, so \textbf{\mr{कुठून}} sits precisely where \textbf{\mr{पुण्याहून}} would sit and nothing moves to the front.

\begin{quote}
\textbf{\mr{तुम्ही} \mr{पुण्याहून} \mr{येता}.} $\to$ \textbf{\mr{तुम्ही} \mr{कुठून} \mr{येता}?}
\end{quote}

Now put the whole family beside each other, because they are one shape with three stems and you have owned two of them since chapter forty:

\begin{quote}
\textbf{\mr{कुठून}} from where · \textbf{\mr{इथून}} from here · \textbf{\mr{तिथून}} from there.
\end{quote}

\subsection*{Your turn}
Say these aloud:

\begin{itemize}
\raggedright
  \item \emph{kuṭhūn} — from where
  \item \emph{tumhī kuṭhūn yetā?} — and answer it about yourself
  \item the three together — \emph{kuṭhūn}, \emph{ithūn}, \emph{tithūn}
\end{itemize}

\begin{tcolorbox}[breakable,colback=teal!4,colframe=teal!35!black,title={Before you move on}]
What two pieces is \textbf{\mr{कुठून}} made of? (\textbf{\mr{कुठे}} plus \textbf{-\mr{हून}}.) Where in the sentence does it stand? (\textbf{Where the answer would stand}.)
\end{tcolorbox}
\section[kuṭhūn? — puṇyāhūn.]{Answer the second interview question}

\label{lesson:MR-R46-from}



\begin{quote}
Say the question and its answer as one exchange, both halves aloud.
\end{quote}

\subsection*{Your turn}
\begin{itemize}
\raggedright
  \item \emph{Say it:} \emph{tumhī kuṭhūn yetā?} — then answer it
  \item \emph{Say it:} \emph{udyāpāsūn} — and hear that this one is a time
  \item \emph{Say it:} \emph{mitrāśī} — and notice this chapter did not touch it
  \item \emph{Recall:} say \textbf{\mr{देणे}} and its root — \emph{deṇe}, from the ancient \emph{deh\textsubscript{3}-}, \textquotedblleft{}to give\textquotedblright{}
\end{itemize}

\begin{tcolorbox}[breakable,colback=teal!4,colframe=teal!35!black,title={Before you move on}]
Which form field can you now turn into a spoken answer? (\textbf{\mr{शहर}} — \textbf{\mr{मी} \mr{पुण्याहून} \mr{येतो}}.) Which of the two endings also does time? (\textbf{\mr{पासून}}.)
\end{tcolorbox}
